\documentclass[10pt,a4paper]{amsart}
\usepackage{amsmath,amssymb,amsthm,graphicx}
\usepackage{verbatim}
\newenvironment{reference}{\comment}{\endcomment}
%\newenvironment{reference}{}{}
\newenvironment{slogan}{\comment}{\endcomment}
\newenvironment{history}{\comment}{\endcomment}

% For commutative diagrams we use Xy-pic
\usepackage[all]{xy}

% We use 2cell for 2-commutative diagrams.
\xyoption{2cell}
% Theorem environments.
%
\theoremstyle{plain}
\newtheorem{theorem}[subsection]{Theorem}
\newtheorem{proposition}[subsection]{Proposition}
\newtheorem{lemma}[subsection]{Lemma}

\theoremstyle{definition}
\newtheorem{definition}[subsection]{Definition}
\newtheorem{example}[subsection]{Example}
\newtheorem{exercise}[subsection]{Exercise}
\newtheorem{situation}[subsection]{Situation}

\theoremstyle{remark}
\newtheorem{remark}[subsection]{Remark}
\newtheorem{remarks}[subsection]{Remarks}

\numberwithin{equation}{subsection}

% Macros
%
\def\lim{\mathop{\mathrm{lim}}\nolimits}
\def\colim{\mathop{\mathrm{colim}}\nolimits}
\def\Spec{\mathop{\mathrm{Spec}}}
\def\Hom{\mathop{\mathrm{Hom}}\nolimits}
\def\Ext{\mathop{\mathrm{Ext}}\nolimits}
\def\SheafHom{\mathop{\mathcal{H}\!\mathit{om}}\nolimits}
\def\SheafExt{\mathop{\mathcal{E}\!\mathit{xt}}\nolimits}
\def\Sch{\mathit{Sch}}
\def\Mor{\mathop{\mathrm{Mor}}\nolimits}
\def\Ob{\mathop{\mathrm{Ob}}\nolimits}
\def\Sh{\mathop{\mathit{Sh}}\nolimits}
\def\NL{\mathop{N\!L}\nolimits}
\def\CH{\mathop{\mathrm{CH}}\nolimits}
\def\proetale{{pro\text{-}\acute{e}tale}}
\def\etale{{\acute{e}tale}}
\def\QCoh{\mathit{QCoh}}
\def\Ker{\mathop{\mathrm{Ker}}}
\def\Im{\mathop{\mathrm{Im}}}
\def\Coker{\mathop{\mathrm{Coker}}}
\def\Coim{\mathop{\mathrm{Coim}}}

% Boxtimes
%
\DeclareMathSymbol{\boxtimes}{\mathbin}{AMSa}{"02}

%
% Macros for moduli stacks/spaces
%
\def\QCohstack{\mathcal{QC}\!\mathit{oh}}
\def\Cohstack{\mathcal{C}\!\mathit{oh}}
\def\Spacesstack{\mathcal{S}\!\mathit{paces}}
\def\Quotfunctor{\mathrm{Quot}}
\def\Hilbfunctor{\mathrm{Hilb}}
\def\Curvesstack{\mathcal{C}\!\mathit{urves}}
\def\Polarizedstack{\mathcal{P}\!\mathit{olarized}}
\def\Complexesstack{\mathcal{C}\!\mathit{omplexes}}
% \Pic is the operator that assigns to X its picard group, usage \Pic(X)
% \Picardstack_{X/B} denotes the Picard stack of X over B
% \Picardfunctor_{X/B} denotes the Picard functor of X over B
\def\Pic{\mathop{\mathrm{Pic}}\nolimits}
\def\Picardstack{\mathcal{P}\!\mathit{ic}}
\def\Picardfunctor{\mathrm{Pic}}
\def\Deformationcategory{\mathcal{D}\!\mathit{ef}}
\usepackage{cleveref}
\crefname{originallemma}{original lemma}{original lemmas}
\Crefname{originallemma}{Original lemma}{Original lemmas}
\makeatletter
\newlabel{lemma-base-change-finite-type}{{03XH}{source}{}{sourcepointer-1}{}}
\newlabel{lemma-base-change-finite-type@cref}{{[originallemma][0][]03XH}{[1][0][]source}{}{}{}}
\makeatother
\crefname{originallemma}{original lemma}{original lemmas}
\Crefname{originallemma}{Original lemma}{Original lemmas}
\makeatletter
\newlabel{lemma-finite-type-points-morphism}{{06EI}{source}{}{sourcepointer-2}{}}
\newlabel{lemma-finite-type-points-morphism@cref}{{[originallemma][0][]06EI}{[1][0][]source}{}{}{}}
\makeatother
\crefname{originallemma}{original lemma}{original lemmas}
\Crefname{originallemma}{Original lemma}{Original lemmas}
\makeatletter
\newlabel{lemma-identify-finite-type-points}{{06EH}{source}{}{sourcepointer-3}{}}
\newlabel{lemma-identify-finite-type-points@cref}{{[originallemma][0][]06EH}{[1][0][]source}{}{}{}}
\makeatother
\title{A locally finite type surjective morphism of algebraic spaces is surjective on finite type points}
\author{The Stacks Project Authors}
\date{stacks-project:spaces-morphisms.tex:local-snapshot-20261012; GNU FDL 1.2 or later}
\begin{document}
\maketitle
\setcounter{subsection}{1}
\noindent\textit{Editorial note.} All delivered mathematical text is the author's own, copied from the pinned Stacks Project chapter file spaces-morphisms.tex; no statement or proof step was rewritten, reordered or removed.\par\medskip
\noindent\textit{Editorial note.} The proof is short but complete: it builds the point on the fibre and transports it back, using the base change stability of local finite typeness in between.\par\medskip
\noindent\textit{Editorial note.} The three references of the counted spans are recorded as non-load-bearing navigation pointers; the argument that consumes them stays inside the counted proof span.\par\medskip
\noindent\textit{Editorial note.} Preamble: the three native excerpt spans of sources/stacks-spaces-morphisms/preamble.tex (lines 8-13, 14-18 and 154-217). The chapter file itself carries no preamble text.\par\medskip
\noindent\textit{Allowed standard prerequisites (editorial).}\begin{itemize}
\item Finite type points of algebraic spaces and their behaviour under locally finite type morphisms. Lemma lemma-finite-type-points-morphism, used at proof lines 5115-5116 and 5126-5128. Gives one inclusion and transports the point found on X\_k back to X.
\item Nonempty algebraic spaces of finite type over a field have a finite type point. Lemma lemma-identify-finite-type-points, used at proof lines 5120-5122. Produces the point x of X\_k.
\item Local finite typeness is stable under base change. Lemma lemma-base-change-finite-type, used at proof lines 5123-5125. Makes X\_k -> X locally of finite type.
\item Representation of a point by a locally finite type morphism from a field. used at proof lines 5117-5118. Supplies the field-valued test point for y.
\end{itemize}\medskip
\ifdefined\hypertarget\hypertarget{sourcepointer-1}{}\fi
\noindent\textit{Source reference.} Original lemma 03XH: focused-seeds/AGP/sources/stacks-snapshot-20261011/tags.txt, source lines 5148--5148. Navigation pointer into the Stacks Project; the standard fact it supplies is applied inside the counted proof span and the argument resting on it stays there.\par\medskip
\ifdefined\hypertarget\hypertarget{sourcepointer-2}{}\fi
\noindent\textit{Source reference.} Original lemma 06EI: focused-seeds/AGP/sources/stacks-snapshot-20261011/tags.txt, source lines 8217--8217. Navigation pointer into the Stacks Project; the standard fact it supplies is applied inside the counted proof span and the argument resting on it stays there.\par\medskip
\ifdefined\hypertarget\hypertarget{sourcepointer-3}{}\fi
\noindent\textit{Source reference.} Original lemma 06EH: focused-seeds/AGP/sources/stacks-snapshot-20261011/tags.txt, source lines 8216--8216. Navigation pointer into the Stacks Project; the standard fact it supplies is applied inside the counted proof span and the argument resting on it stays there.\par\medskip
\begin{lemma}
\label{lemma-finite-type-points-surjective-morphism}
Let $S$ be a scheme.
Let $f : X \to Y$ be a morphism of algebraic spaces over $S$.
If $f$ is locally of finite type and surjective, then
$f(X_{\text{ft-pts}}) = Y_{\text{ft-pts}}$.
\end{lemma}

\begin{proof}
We have $f(X_{\text{ft-pts}}) \subset Y_{\text{ft-pts}}$ by
Lemma \ref{lemma-finite-type-points-morphism}.
Let $y \in |Y|$ be a finite type point. Represent $y$ by a morphism
$\Spec(k) \to Y$ which is locally of finite type.
As $f$ is surjective the algebraic space
$X_k = \Spec(k) \times_Y X$ is nonempty, therefore
has a finite type point $x \in |X_k|$ by
Lemma \ref{lemma-identify-finite-type-points}.
Now $X_k \to X$ is a morphism which is locally of finite type as a base change
of $\Spec(k) \to Y$
(Lemma \ref{lemma-base-change-finite-type}).
Hence the image of $x$ in $X$ is a finite type point by
Lemma \ref{lemma-finite-type-points-morphism}
which maps to $y$ by construction.
\end{proof}

\end{document}
